\documentclass[12pt,a4paper]{article}
\usepackage[utf8]{inputenc}
\usepackage{geometry}
\geometry{a4paper, margin=1in}
\usepackage{graphicx}
\usepackage{hyperref}
\usepackage{listings}
\usepackage{xcolor}
\usepackage{booktabs}

\lstdefinestyle{jsonstyle}{
    basicstyle=\ttfamily\footnotesize,
    keywordstyle=\color{blue},
    stringstyle=\color{teal},
    commentstyle=\color{gray},
    showstringspaces=false,
    breaklines=true,
    frame=single,
    backgroundcolor=\color{black!5}
}

\title{\textbf{Athena Surveillance System: Camera Control API Integration Report}}
\author{Technical Documentation}
\date{\today}

\begin{document}
\maketitle

\begin{abstract}
This report provides an in-depth analysis of the Camera Control API utilized within the Athena Surveillance System. Moving away from standard, high-latency XML/SOAP-based ONVIF protocols, Athena implements a highly optimized, low-latency integration with the native Dahua/Hikvision \textbf{RPC2 HTTP API}. This document outlines the API's authentication protocols, heartbeat session management, PTZ action dispatching, and auxiliary configuration controls.
\end{abstract}

\section{Overview of the RPC2 Interface}
The \texttt{RPC2} (Remote Procedure Call) interface is a JSON-based protocol exposed over standard HTTP. It allows direct manipulation of the camera's internal operating system. The API communicates over two primary endpoints:
\begin{itemize}
    \item \texttt{http://<ip>:<port>/RPC2\_Login}: Used strictly for session negotiation and cryptographic authentication.
    \item \textbf{\texttt{http://<ip>:<port>/RPC2}}: The main command endpoint for authenticated sessions.
\end{itemize}
All requests are structured as JSON payloads containing a \texttt{method}, \texttt{params}, a request \texttt{id}, and a \texttt{session} identifier.

\section{Authentication Protocol}
Authentication is governed by the \texttt{RPCSessionPool} class, which executes a two-phase, challenge-response mechanism using MD5 hashing. This prevents plaintext credential transmission.

\subsection{Phase 1: Challenge Request}
The server initializes the login sequence by requesting a cryptographic challenge from the camera.
\begin{lstlisting}[style=jsonstyle, language=json, caption=Phase 1 Request]
{
    "method": "global.login",
    "params": {
        "userName": "admin",
        "password": "plaintext_password",
        "clientType": "Web3.0"
    },
    "id": 1,
    "session": 0
}
\end{lstlisting}
The camera rejects the plaintext password but responds with a \texttt{realm}, a \texttt{random} challenge string, and a temporary \texttt{session} ID.

\subsection{Phase 2: Cryptographic Response}
The backend computes a nested MD5 hash using the received parameters:
\begin{enumerate}
    \item $h_1 = \text{MD5}(\text{username} : \text{realm} : \text{password})$
    \item $h_2 = \text{MD5}(\text{username} : \text{random} : h_1)$
\end{enumerate}

The hashed response ($h_2$) is then submitted back to the \texttt{RPC2\_Login} endpoint:
\begin{lstlisting}[style=jsonstyle, language=json, caption=Phase 2 Request]
{
    "method": "global.login",
    "params": {
        "userName": "admin",
        "password": "<h2_hash_value>",
        "clientType": "Web3.0"
    },
    "id": 2,
    "session": "<temporary_session_id>"
}
\end{lstlisting}
Upon success, the camera validates the hash and solidifies the \texttt{session\_id} for subsequent RPC calls.

\section{Session Management \& Keep-Alive}
To eliminate the latency of re-authenticating before every PTZ movement, the \texttt{RPCSessionPool} maintains a persistent session cache in memory. 

To prevent the camera from terminating idle sessions, a background thread (\texttt{RPC-Heartbeat}) pings the \texttt{RPC2} endpoint every 30 seconds using a lightweight diagnostic method:
\begin{lstlisting}[style=jsonstyle, language=json, caption=Heartbeat Request]
{
    "method": "magicBox.getDeviceType",
    "params": null,
    "id": 1,
    "session": "<active_session_id>"
}
\end{lstlisting}
If the heartbeat fails or the session expires, the pool automatically invalidates the cache and triggers a transparent re-authentication.

\section{PTZ Action Dispatching}
The core function of the API integration is high-speed Pan-Tilt-Zoom (PTZ) control, driven by the AI tracking engine's PID controllers.

\subsection{Movement Commands}
Movement is triggered using the \texttt{ptz.start} and \texttt{ptz.stop} methods. The \texttt{code} parameter dictates the axis (\texttt{Up}, \texttt{Down}, \texttt{Left}, \texttt{Right}), while \texttt{arg1} and \texttt{arg2} map to the calculated speed output (ranging from 1 to 8).

\begin{lstlisting}[style=jsonstyle, language=json, caption=PTZ Start Command]
{
    "method": "ptz.start",
    "params": {
        "code": "Right",
        "arg1": 5,
        "arg2": 5,
        "arg3": 0,
        "arg4": 0
    },
    "object": 1, // Camera Channel
    "id": 388,
    "session": "<active_session_id>"
}
\end{lstlisting}

When the target centers in the frame, the engine issues a \texttt{ptz.stop} command for the respective axis to halt the motors.

\section{Auxiliary Controls (Lighting \& Telemetry)}
Beyond PTZ, the RPC API accesses deeper configuration layers of the camera hardware.

\subsection{Smart Illumination}
The system can toggle the camera's physical LEDs or IR illuminators via the \texttt{configManager.setConfig} method.
\begin{lstlisting}[style=jsonstyle, language=json, caption=Lighting Configuration]
{
    "method": "configManager.setConfig",
    "params": {
        "name": "Lighting",
        "table": [[{"Mode": "Auto"}]]
    },
    "id": 389,
    "session": "<active_session_id>"
}
\end{lstlisting}

\subsection{Virtual Telemetry via Dead-Reckoning}
Because native polling of PTZ absolute coordinates over HTTP is too slow for 30 FPS feedback loops, the API control layer utilizes a \textbf{Dead-Reckoning} algorithm. The backend integrates the time-delta of active \texttt{ptz.start} commands multiplied by their speed vectors to estimate the physical pan and tilt angles in real-time, completely avoiding network-bound telemetry bottlenecks.

\section{Conclusion}
The Athena Surveillance System's migration to the native Dahua RPC2 JSON API represents a crucial optimization. By implementing persistent cryptographic session pooling, lightweight heartbeat pinging, and JSON-based PTZ dispatching, the platform achieves the microsecond responsiveness strictly required for closed-loop AI-driven robotic camera tracking.

\end{document}
